\begin{frame}{FAQ (2/2): $\gamma$ 의 극단}
  \begin{itemize}
    \item \textbf{Q. $\gamma = 0$ 는 무슨 뜻인가요?}
      \begin{itemize}
        \item ``당장의 스텝 \textbf{하나}의 보상뿐''을 본다는 뜻.
          $G_t = R_t$ 로 리턴이 즉시 보상과 같아지고,
          \textbf{근시안의 극단}.
        \item MDP는 퇴화하여 \kb{다중팔 밴딧}(각 상태가 별개의
          밴딧)이 된다.
        \item 드물지만 ``보상이 이미 \textbf{누적된 값}으로
          주어지는 환경''에서는 $\gamma = 0$ 으로 ``최종
          총점''을 한 번만 세는 것이 \textbf{이중 계산을 막는
          정석}.
      \end{itemize}
    \item \textbf{Q. $\gamma$ 를 0.999처럼 1에 아주 가깝게 하면
      ``더 좋은'' 것 아닌가?}
      \begin{itemize}
        \item (1) \textbf{수학}: $\gamma = 1$ 이면 끝나지 않는
          과업에서 발산할 수 있다. $\gamma < 1$ 은 \textbf{수렴을
          수학적으로 보장}.
        \item (2) \textbf{실용}: 1에 가까울수록 리턴의
          \textbf{분산}이 커지고, 수렴이 \textbf{느려져} 학습이
          불안정해질 수 있다.
        \item ``1에 가까울수록 좋은가''는 \textbf{환경의 수명}과
          \textbf{알고리즘의 분산}을 함께 저울질해야 정해지는
          \textbf{설계 판단}.
      \end{itemize}
  \end{itemize}
\end{frame}
